% CCF-1688 源宝合作基金项目申报书专属稳定样式。
% 项目正文与表格数据保留在 projects/ccf-1688/extraTex/。

\RequirePackage[table]{xcolor}
\RequirePackage{geometry}
\RequirePackage{graphicx}
\RequirePackage{fontspec}
\RequirePackage{xeCJK}
\RequirePackage{array}
\RequirePackage{fancyhdr}
\RequirePackage{amsmath}
\RequirePackage{titlesec}
\RequirePackage{indentfirst}
\RequirePackage{caption}
\DeclareCaptionFont{ccftablecaption}{\fontsize{10pt}{11.5pt}\selectfont}

\expandafter\geometry\expandafter{\NSFCGeometryOptions}
\setlength{\headheight}{15pt}
\setlength{\headsep}{17.5pt}
\setlength{\footskip}{36pt}
\raggedbottom

\definecolor{CCFBlue}{RGB}{0,102,255}
\definecolor{CCFHintGray}{RGB}{127,127,127}

% CCF 正文中的公式、参考文献 URL 与交叉引用均需可正常渲染和跳转。
\RequirePackage{hyperref}
\RequirePackage{xurl}
\hypersetup{%
  colorlinks=true,
  linkcolor=black,
  citecolor=black,
  urlcolor=CCFBlue,
  CJKbookmarks=true
}
\urlstyle{same}

% 使用共享字体资源：中文保持 SimSun / Noto Sans SC，英文与数字统一使用
% Times New Roman，保证正文、标题、表格、页眉和参考文献的西文字体一致。
\newfontfamily\ccfsonglatinfont[
  Path={\NSFCAssetFontsDir},
  UprightFont=TimesNewRoman.ttf,
  BoldFont=TimesNewRoman.ttf,
  BoldFeatures={FakeBold=2.5},
  ItalicFont=TimesNewRoman.ttf,
  ItalicFeatures={FakeSlant=0.2},
  BoldItalicFont=TimesNewRoman.ttf,
  BoldItalicFeatures={FakeBold=2.5,FakeSlant=0.2}
]{TimesNewRoman.ttf}
\setCJKfamilyfont{ccfsong}[
  Path={\NSFCAssetFontsDir},
  AutoFakeBold=2.5,
  AutoFakeSlant=0.2
]{SimSun.ttf}
\newfontfamily\ccfsymbolfont[
  Path={\NSFCAssetFontsDir},
  UprightFont=SimSun.ttf,
  BoldFont=SimSun.ttf,
  BoldFeatures={FakeBold=2.5}
]{SimSun.ttf}

\newfontfamily\ccfsanslatinfont[
  Path={\NSFCAssetFontsDir},
  UprightFont=TimesNewRoman.ttf,
  BoldFont=TimesNewRoman.ttf,
  BoldFeatures={FakeBold=2.5},
  ItalicFont=TimesNewRoman.ttf,
  ItalicFeatures={FakeSlant=0.2},
  BoldItalicFont=TimesNewRoman.ttf,
  BoldItalicFeatures={FakeBold=2.5,FakeSlant=0.2}
]{TimesNewRoman.ttf}
\setCJKfamilyfont{ccfsans}[
  Path={\NSFCAssetFontsDir},
  BoldFont=NotoSansSC-Bold.otf,
  AutoFakeSlant=0.2
]{NotoSansSC-Regular.otf}

\newcommand{\ccfsonglatin}{\ccfsonglatinfont}
\newcommand{\ccfsong}{\CJKfamily{ccfsong}}
\newcommand{\ccfnsonglatin}{\ccfsonglatinfont}
\newcommand{\ccfnsong}{\CJKfamily{ccfsong}}
\newcommand{\ccfsanslatin}{\ccfsanslatinfont}
\newcommand{\ccfsans}{\CJKfamily{ccfsans}}
\newcommand{\CCFSymbol}[1]{{\ccfsymbolfont #1}}

% 参考文献沿用国自然模板的紧凑列表口径：五号字、单倍行距、条目间距为 0，
% 并保留悬挂编号与两端对齐；标题仍使用 CCF 中文字体和黑色标题样式。
\newlength{\CCFBibTitleAboveSkip}
\newlength{\CCFBibTitleBelowSkip}
\newlength{\CCFBibItemSep}
\setlength{\CCFBibTitleAboveSkip}{10pt}
\setlength{\CCFBibTitleBelowSkip}{10pt}
\setlength{\CCFBibItemSep}{0pt}

\makeatletter
\renewenvironment{thebibliography}[1]{%
  \par\addvspace{\CCFBibTitleAboveSkip}%
  \noindent{\ccfsonglatin\ccfsong\fontsize{14.1pt}{17pt}\selectfont\bfseries\refname}\par
  \vspace{\CCFBibTitleBelowSkip}%
  \ccfsonglatin\ccfsong\fontsize{10.5pt}{14.45pt}\selectfont
  \frenchspacing
  \list{\@biblabel{\@arabic\c@enumiv}}{%
    \settowidth\labelwidth{\@biblabel{#1}}%
    \leftmargin\labelwidth
    \advance\leftmargin\labelsep
    \setlength{\rightmargin}{0pt}%
    \setlength{\itemsep}{\CCFBibItemSep}%
    \setlength{\parsep}{0pt}%
    \setlength{\parskip}{0pt}%
    \setlength{\topsep}{0pt}%
    \setlength{\partopsep}{0pt}%
    \usecounter{enumiv}%
    \let\p@enumiv\@empty
    \renewcommand\theenumiv{\@arabic\c@enumiv}%
  }%
  \spaceskip=0.25em plus 0.25em minus 0.15em\relax
  \xspaceskip=\spaceskip
  \setlength{\emergencystretch}{0pt}%
  \tolerance=2000\relax
  \clubpenalty4000\widowpenalty4000%
  \sfcode`\.=1000\relax
}{%
  \def\@noitemerr{\@latex@warning{Empty `thebibliography' environment}}%
  \endlist
}
\makeatother

\setCJKmainfont[
  Path={\NSFCAssetFontsDir},
  AutoFakeBold=2.5,
  AutoFakeSlant=0.2
]{SimSun.ttf}
\setCJKmonofont[Path={\NSFCAssetFontsDir}]{SimSun.ttf}
\setmainfont[
  Path={\NSFCAssetFontsDir},
  UprightFont=TimesNewRoman.ttf,
  BoldFont=TimesNewRoman.ttf,
  BoldFeatures={FakeBold=2.5},
  ItalicFont=TimesNewRoman.ttf,
  ItalicFeatures={FakeSlant=0.2},
  BoldItalicFont=TimesNewRoman.ttf,
  BoldItalicFeatures={FakeBold=2.5,FakeSlant=0.2}
]{TimesNewRoman.ttf}
\setmonofont[
  Path={\NSFCAssetFontsDir},
  UprightFont=TimesNewRoman.ttf,
  BoldFont=TimesNewRoman.ttf,
  BoldFeatures={FakeBold=2.5},
  ItalicFont=TimesNewRoman.ttf,
  ItalicFeatures={FakeSlant=0.2},
  BoldItalicFont=TimesNewRoman.ttf,
  BoldItalicFeatures={FakeBold=2.5,FakeSlant=0.2}
]{TimesNewRoman.ttf}

\AtBeginDocument{%
  \renewcommand{\baselinestretch}{1.0}%
  \fontsize{12.05pt}{23.4pt}\selectfont\ccfsonglatin\ccfsong
  \setlength{\parindent}{2em}%
  \setlength{\parskip}{0pt}%
  \setlength{\emergencystretch}{3em}%
  % 表题采用国自然申请书的紧凑层级：小于正文、居中，并与表体留白。
  % 对置于表体上方的 caption，skip 控制表题与表体之间的距离。
  \captionsetup[table]{font=ccftablecaption,skip=6pt}%
}
\XeTeXlinebreaklocale "zh"

\fancypagestyle{ccf1688}{%
  \fancyhf{}%
  \fancyhead[C]{\raisebox{1.4pt}{{\ccfsanslatin\ccfsans\fontsize{12.05pt}{15pt}\selectfont\bfseries CONFIDENTIAL}}}%
  \fancyfoot[C]{{\ccfsanslatin\ccfsans\fontsize{9pt}{11pt}\selectfont 保密信息}}%
  \renewcommand{\headrulewidth}{0.75pt}%
  \renewcommand{\footrulewidth}{0pt}%
}
\pagestyle{ccf1688}

\newcolumntype{C}[1]{>{\centering\arraybackslash}m{#1}}
\newcolumntype{L}[1]{>{\raggedright\arraybackslash}m{#1}}
\newcommand{\CCFCenteredFixedCell}[2]{%
  \parbox[c][#1][c]{\linewidth}{\centering #2\par}%
}
\newcommand{\CCFLeftFixedCell}[2]{%
  \parbox[c][#1][c]{\linewidth}{\raggedright #2\par}%
}

\newcommand{\CCFHint}[1]{{\fontsize{11.09pt}{23.4pt}\selectfont\color{CCFHintGray}#1}}
\newcommand{\CCFBlueHeader}{\rowcolor{CCFBlue}}
\newcommand{\CCFTableHeader}[1]{{\ccfsonglatin\ccfsong\fontsize{12.05pt}{14.5pt}\selectfont\bfseries\color{white}#1}}
\newcommand{\CCFNormalTableHeader}[1]{\CCFTableHeader{#1}}
\newcommand{\CCFSmallTableHeader}[1]{{\ccfsonglatin\ccfsong\fontsize{10pt}{12pt}\selectfont\bfseries\color{white}#1}}
\newcommand{\CCFTableSongHint}[1]{{\ccfsonglatin\ccfsong\fontsize{11.09pt}{13pt}\selectfont\color{CCFHintGray}#1}}
\newcommand{\CCFTableSansHint}[1]{{\ccfsanslatin\ccfsans\fontsize{10.5pt}{13pt}\selectfont\itshape\color{CCFHintGray}#1}}

\NewDocumentCommand{\CCFMainHeading}{O{24pt} m m}{%
  \noindent\makebox[24pt][l]{\bfseries #2}{\bfseries #3}\par
}

% 二级标题由 CCF 官方骨架手工编号；此宏在排标题时同步 LaTeX 计数器，
% 使其下的三级标题能得到 1.2.1、1.4.1 等正确层级编号。
\NewDocumentCommand{\CCFSecondHeading}{m m m}{%
  \setcounter{section}{#1}%
  \setcounter{subsection}{#2}%
  \setcounter{subsubsection}{0}%
  \noindent #1.#2.\hspace{6pt}#3\par
}

\setcounter{secnumdepth}{4}
\renewcommand{\thesubsubsection}{\thesection.\arabic{subsection}.\arabic{subsubsection}}
\titleformat{\subsubsection}
  {\ccfsonglatin\ccfsong\fontsize{12.05pt}{18pt}\selectfont\bfseries}
  {\thesubsubsection}
  {0.5em}
  {}
\titlespacing*{\subsubsection}{22pt}{6pt}{0pt}

\titleclass{\subsubsubsection}{straight}[\subsubsection]
\newcounter{subsubsubsection}[subsubsection]
\renewcommand{\thesubsubsubsection}{（\arabic{subsubsubsection}）}
\titleformat{\subsubsubsection}
  {\ccfsonglatin\ccfsong\fontsize{12.05pt}{18pt}\selectfont\bfseries}
  {\thesubsubsubsection}
  {0.25em}
  {}
\titlespacing*{\subsubsubsection}{22pt}{4pt}{0pt}
\makeatletter
\def\toclevel@subsubsubsection{4}
\def\l@subsubsubsection{\@dottedtocline{4}{7em}{4em}}
\makeatother

\newcommand{\CCFCoverInfoTable}{%
  \begingroup
  \setlength{\tabcolsep}{6pt}%
  \setlength{\arrayrulewidth}{0.5pt}%
  \fontsize{10.5pt}{13pt}\selectfont
  \noindent
  \begin{tabular}{|C{113.3pt}|C{278.04pt}|}
    \hline
    \CCFCenteredFixedCell{31.5pt}{\bfseries 申请者姓名} & \CCFCenteredFixedCell{31.5pt}{\CCFApplicantName} \\ \hline
    \CCFCenteredFixedCell{31.5pt}{\bfseries CCF会员码} & \CCFCenteredFixedCell{31.5pt}{\CCFMemberCode} \\ \hline
    \CCFCenteredFixedCell{31.5pt}{\bfseries 所属学校/科研机构} &
      \CCFCenteredFixedCell{31.5pt}{\CCFInstitution\ifx\CCFInstitution\empty{\color{CCFHintGray}（提示：申报单位与项目科研协议签署单位需保持一致）}\fi} \\ \hline
    \CCFCenteredFixedCell{31.5pt}{\bfseries 提交日期} & \CCFCenteredFixedCell{31.5pt}{\CCFSubmissionDate} \\ \hline
  \end{tabular}%
  \par
  \vspace{-4.5pt}%
  \endgroup
}
